\documentclass[10pt,a4paper]{amsart}
\usepackage{amsmath,amssymb,amsthm,graphicx}
% Retained preamble for the Stacks Project chapter morphisms.tex (Morphisms of Schemes).
% The chapter file pinned-original-source.tex carries only \input{preamble} on its first line, so no
% preamble text lives inside the chapter itself.  The Stacks Project uses one single shared preamble for
% every chapter; the snapshot stages it verbatim in sources/stacks-chow/preamble.tex.  The block below is
% copied mechanically, line for line, from that file's lines 14-18 (the Xy-pic setup) and lines 154-217
% (the native theorem declarations together with the mathematical macros).  Nothing was added, removed or
% redefined.  The wrapper class and the amsmath/amssymb/amsthm/graphicx base packages are supplied by the
% assembler, so the class/\externaldocument part of the shared preamble is not repeated here.
% For commutative diagrams we use Xy-pic
\usepackage[all]{xy}

% We use 2cell for 2-commutative diagrams.
\xyoption{2cell}
% Theorem environments.
%
\theoremstyle{plain}
\newtheorem{theorem}[subsection]{Theorem}
\newtheorem{proposition}[subsection]{Proposition}
\newtheorem{lemma}[subsection]{Lemma}

\theoremstyle{definition}
\newtheorem{definition}[subsection]{Definition}
\newtheorem{example}[subsection]{Example}
\newtheorem{exercise}[subsection]{Exercise}
\newtheorem{situation}[subsection]{Situation}

\theoremstyle{remark}
\newtheorem{remark}[subsection]{Remark}
\newtheorem{remarks}[subsection]{Remarks}

\numberwithin{equation}{subsection}

% Macros
%
\def\lim{\mathop{\mathrm{lim}}\nolimits}
\def\colim{\mathop{\mathrm{colim}}\nolimits}
\def\Spec{\mathop{\mathrm{Spec}}}
\def\Hom{\mathop{\mathrm{Hom}}\nolimits}
\def\Ext{\mathop{\mathrm{Ext}}\nolimits}
\def\SheafHom{\mathop{\mathcal{H}\!\mathit{om}}\nolimits}
\def\SheafExt{\mathop{\mathcal{E}\!\mathit{xt}}\nolimits}
\def\Sch{\mathit{Sch}}
\def\Mor{\mathop{\mathrm{Mor}}\nolimits}
\def\Ob{\mathop{\mathrm{Ob}}\nolimits}
\def\Sh{\mathop{\mathit{Sh}}\nolimits}
\def\NL{\mathop{N\!L}\nolimits}
\def\CH{\mathop{\mathrm{CH}}\nolimits}
\def\proetale{{pro\text{-}\acute{e}tale}}
\def\etale{{\acute{e}tale}}
\def\QCoh{\mathit{QCoh}}
\def\Ker{\mathop{\mathrm{Ker}}}
\def\Im{\mathop{\mathrm{Im}}}
\def\Coker{\mathop{\mathrm{Coker}}}
\def\Coim{\mathop{\mathrm{Coim}}}

% Boxtimes
%
\DeclareMathSymbol{\boxtimes}{\mathbin}{AMSa}{"02}

%
% Macros for moduli stacks/spaces
%
\def\QCohstack{\mathcal{QC}\!\mathit{oh}}
\def\Cohstack{\mathcal{C}\!\mathit{oh}}
\def\Spacesstack{\mathcal{S}\!\mathit{paces}}
\def\Quotfunctor{\mathrm{Quot}}
\def\Hilbfunctor{\mathrm{Hilb}}
\def\Curvesstack{\mathcal{C}\!\mathit{urves}}
\def\Polarizedstack{\mathcal{P}\!\mathit{olarized}}
\def\Complexesstack{\mathcal{C}\!\mathit{omplexes}}
% \Pic is the operator that assigns to X its picard group, usage \Pic(X)
% \Picardstack_{X/B} denotes the Picard stack of X over B
% \Picardfunctor_{X/B} denotes the Picard functor of X over B
\def\Pic{\mathop{\mathrm{Pic}}\nolimits}
\def\Picardstack{\mathcal{P}\!\mathit{ic}}
\def\Picardfunctor{\mathrm{Pic}}
\def\Deformationcategory{\mathcal{D}\!\mathit{ef}}
\usepackage{cleveref}
\crefname{originallemma}{original lemma}{original lemmas}
\Crefname{originallemma}{Original lemma}{Original lemmas}
\makeatletter
\newlabel{lemma-affine-separated}{{lemma-affine-separated}{source}{}{sourcepointer-1}{}}
\newlabel{lemma-affine-separated@cref}{{[originallemma][0][]lemma-affine-separated}{[1][0][]source}{}{}{}}
\makeatother
\crefname{originallemma}{original lemma}{original lemmas}
\Crefname{originallemma}{Original lemma}{Original lemmas}
\makeatletter
\newlabel{lemma-quasi-coherence-scalar-restriction}{{lemma-quasi-coherence-scalar-restriction}{source}{}{sourcepointer-2}{}}
\newlabel{lemma-quasi-coherence-scalar-restriction@cref}{{[originallemma][0][]lemma-quasi-coherence-scalar-restriction}{[1][0][]source}{}{}{}}
\makeatother
\crefname{originalreference}{original reference}{original references}
\Crefname{originalreference}{Original reference}{Original references}
\makeatletter
\newlabel{schemes-lemma-push-forward-quasi-coherent}{{schemes-lemma-push-forward-quasi-coherent}{source}{}{sourcepointer-3}{}}
\newlabel{schemes-lemma-push-forward-quasi-coherent@cref}{{[originalreference][0][]schemes-lemma-push-forward-quasi-coherent}{[1][0][]source}{}{}{}}
\makeatother
\title{For an affine morphism the pushforward is an equivalence from quasi-coherent modules on the source to quasi-coherent modules over the pushforward algebra}
\author{The Stacks Project Authors}
\date{stacks-project:morphisms.tex:local-snapshot-20261012; GNU FDL 1.2 or later}
\begin{document}
\maketitle
\noindent\textit{Editorial note.} All delivered mathematical text is the author's own, copied from the pinned Stacks Project chapter file pinned-original-source.tex; no statement or proof step was rewritten, reordered or removed.\par\medskip
\noindent\textit{Editorial note.} Following the multi-claim convention for this source, source\_url is the source-level URL https://stacks.math.columbia.edu/ and the source label lemma-affine-equivalence-modules is carried in claim\_key and here rather than in a per-claim tag URL.\par\medskip
\noindent\textit{Editorial note.} The preamble is the chapter's shared Stacks Project preamble, retained as sources/stacks-morphisms/retained-preamble.tex: a line-for-line excerpt of the shared Stacks preamble staged with this snapshot. No macro was invented or redefined, and this card adds no preamble text of its own.\par\medskip
\noindent\textit{Editorial note.} Statement span 1895-1919 and proof span 1921-1967 come from the halfA pool card operations/pool-cards/stacks-morphisms-halfA.json and were re-read in the pinned chapter file before this card was written.\par\medskip
\noindent\textit{Editorial note.} The author's first sentence, 'The final statement is Lemma [quasi-coherence of scalar restriction].', is a citation rather than a proof, and is preserved verbatim. The equivalence of categories, which is the main assertion, is proved in the counted span.\par\medskip
\noindent\textit{Editorial note.} The proof ends with the author's sentence 'We omit the details.' after stating the key algebra identity (N tensor A) tensor over (A tensor A) with A equals N; that sentence is kept verbatim and what it leaves to the reader is the verification of that identity in the affine case, the identity itself being displayed in the counted span.\par\medskip
\noindent\textit{Editorial note.} 3 cross-references occur inside the counted spans; they are listed in external\_references as navigation pointers only.\par\medskip
\noindent\textit{Allowed standard prerequisites (editorial).}\begin{itemize}
\item An affine morphism is separated, and the pushforward of a quasi-coherent module along it is a quasi-coherent module on the base, cited by the author from the same chapter and from the schemes chapter. Cited at proof lines 1922-1925. Makes the pushforward land in quasi-coherent sheaves on the base, hence defines the functor of the statement.
\item Adjunction between pullback and pushforward, and the algebra structure on the pullback of the pushforward algebra together with its map to the structure sheaf. Used at proof lines 1928-1937 to define the quasi-inverse and at 1938-1945 to build the two comparison maps. Supplies the quasi-inverse as a tensor product over the pulled-back algebra and the unit and counit transformations.
\item The same-chapter statement that a module over a quasi-coherent algebra is quasi-coherent over the base if and only if it is quasi-coherent over the algebra. This is the author's first sentence, at proof line 1921, disposing of the final assertion by citation. It is the author's own disposition of the moreover part of the statement; the counted proof itself constructs the equivalence of categories.
\end{itemize}\medskip
\ifdefined\hypertarget\hypertarget{sourcepointer-1}{}\fi
\noindent\textit{Source reference.} Original lemma lemma-affine-separated: sources/stacks-morphisms/pinned-original-source.tex, source lines 1721--1721. Same-chapter reference used in the retained author text at source line(s) 1923, 11481. It is supplied only as a navigation pointer so that the retained text compiles; the counted item delivers the author's own statement and the author's complete proof, and no part of the referenced item is copied in or substituted for the counted argument.\par\medskip
\ifdefined\hypertarget\hypertarget{sourcepointer-2}{}\fi
\noindent\textit{Source reference.} Original lemma lemma-quasi-coherence-scalar-restriction: sources/stacks-morphisms/pinned-original-source.tex, source lines 1845--1845. Same-chapter reference used in the retained author text at source line(s) 1922. It is supplied only as a navigation pointer so that the retained text compiles; the counted item delivers the author's own statement and the author's complete proof, and no part of the referenced item is copied in or substituted for the counted argument.\par\medskip
\ifdefined\hypertarget\hypertarget{sourcepointer-3}{}\fi
\noindent\textit{Source reference.} Original reference schemes-lemma-push-forward-quasi-coherent: sources/stacks-morphisms/pinned-original-source.tex, source lines 262--262. Same-chapter reference used in the retained author text at source line(s) 262, 404, 882, 1765, 1924, 2256, 2296, 8726, 9175, 10731, 13903. It is supplied only as a navigation pointer so that the retained text compiles; the counted item delivers the author's own statement and the author's complete proof, and no part of the referenced item is copied in or substituted for the counted argument.\par\medskip
\begin{lemma}
\label{lemma-affine-equivalence-modules}
Let $f : X \to S$ be an affine morphism of schemes.
Let $\mathcal{A} = f_*\mathcal{O}_X$.
The functor $\mathcal{F} \mapsto f_*\mathcal{F}$ induces
an equivalence of categories
$$
\left\{
\begin{matrix}
\text{category of quasi-coherent}\\
\mathcal{O}_X\text{-modules}
\end{matrix}
\right\}
\longrightarrow
\left\{
\begin{matrix}
\text{category of quasi-coherent}\\
\mathcal{A}\text{-modules}
\end{matrix}
\right\}
$$
Moreover, an $\mathcal{A}$-module is
quasi-coherent as an $\mathcal{O}_S$-module if and only if
it is quasi-coherent as an $\mathcal{A}$-module.
\end{lemma}

\begin{proof}
The final statement is Lemma \ref{lemma-quasi-coherence-scalar-restriction}.
By Lemma \ref{lemma-affine-separated} and
Schemes, Lemma \ref{schemes-lemma-push-forward-quasi-coherent}
the pushforward $f_*\mathcal{F}$ of a quasi-coherent $\mathcal{O}_X$-module
$\mathcal{F}$
is a quasi-coherent $\mathcal{O}_S$-module. Hence a functor as in
the statement of the lemma.

\medskip\noindent
We will construct an quasi-inverse $h$ to this functor. Let $\mathcal{G}$
be a quasi-coherent $\mathcal{A}$-module. Then we set
$$
h(\mathcal{G}) = f^*\mathcal{G} \otimes_{f^*\mathcal{A}} \mathcal{O}_X =
f^*\mathcal{G} \otimes_{f^*f_*\mathcal{O}_X} \mathcal{O}_X
$$
Elucidation: the pullback $f^*\mathcal{A} = f^*f_*\mathcal{O}_X$
is an $\mathcal{O}_X$-algebra, the adjunction map
$f^*f_*\mathcal{O}_X \to \mathcal{O}_X$ is an algebra homomorphism, and
the pullback $f^*\mathcal{G}$ is an $f^*\mathcal{A}$-module. Observe that
$h(\mathcal{G})$ is quasi-coherent as quasi-coherence is preserved by
pullbacks and change of rings. Observe that there is a functorial map
$$
h(f_*\mathcal{F}) =
f^*f_*\mathcal{F} \otimes_{f^*f_*\mathcal{O}_X} \mathcal{O}_X \to \mathcal{F}
$$
coming from the adjunction map $f^*f_*\mathcal{F} \to \mathcal{F}$ and
a functorial map
$$
\mathcal{G} \to f_*h(\mathcal{G}) \quad\text{adjoint to the map}\quad
f^*\mathcal{G} \to f^*\mathcal{G} \otimes_{f^*f_*\mathcal{O}_X} \mathcal{O}_X
$$
which sends a local section $s$ of $f^*\mathcal{F}$ to $s \otimes 1$.
To finish the proof it suffices to show that these maps are isomorphisms
for $\mathcal{F}$ and $\mathcal{G}$ as above. This may be checked on the
members of an affine covering, i.e., when $X$ and $S$ are affine.

\medskip\noindent
The key algebra observation which makes this work is the
following: Let $R \to A$ be a ring map. Let $N$ be an $A$-module. Then
$$
(N \otimes_R A) \otimes_{(A \otimes_R A)} A = N
$$
Namely, the left hand side of this equality is the effect of applying
$h$ to the quasi-coherent $\widetilde{A}$-module $\widetilde{N}$ on $\Spec(R)$.
We omit the details.
\end{proof}

\end{document}
